\section{Governing equations}
\label{sec:governing}

We consider an incompressible medium on a fixed domain $\Omega\subset\mathbb{R}^2$
occupied by a Newtonian fluid and one or more finite-strain solid bodies. A single
velocity field $\uv(\xv,t)$ and pressure $p(\xv,t)$ are shared by both phases. Momentum
and incompressibility read
\begin{align}
\rho\left(\partial_t\uv + (\uv\!\cdot\!\grad)\uv\right)
   &= -\grad p + \dive\sig_{\mathrm{d}} + \bm f_{\Gamma} + \bm f_{c},
   \label{eq:mom}\\
\dive\uv &= 0,
   \label{eq:incomp}
\end{align}
where $\rho$ is the (piecewise-constant) density, $\bm f_\Gamma$ is the surface-tension
force and $\bm f_c$ the contact force. The deviatoric stress $\sig_{\mathrm d}$ is
blended between fluid and solid through a smoothed indicator (one-fluid model),
\begin{equation}
\sig_{\mathrm d} = 2\mu_f\,\bm D
   + (1-H)\,\sig_{\!s},
   \qquad
\bm D=\tfrac12\!\left(\grad\uv+\grad\uv^{\!\top}\right),
   \label{eq:blend}
\end{equation}
with $\mu_f$ the fluid viscosity, $\bm D$ the rate-of-strain, $\sig_{\!s}$ the solid
Cauchy stress (\cref{sec:rmt,sec:constitutive}), and $H\in[0,1]$ a smoothed Heaviside
of the level set $\phi$ ($H\to1$ in the fluid, $H\to0$ in the solid). The trace of
$\sig_{\!s}$ is carried by the pressure, so only its deviatoric part enters
\cref{eq:mom} up to a redefinition of $p$.

The material occupied by the solid is tracked by the reference map (\cref{sec:rmt}),
from which the level set $\phi$ and the deformation are reconstructed. Boundary
conditions are problem-dependent (no-slip walls with an optional moving lid, free-slip
impermeable walls, or periodic); the pressure boundary condition is chosen consistently
with the velocity boundary condition (\cref{sec:discretization}).
